\documentclass[main.tex]{subfiles}


\begin{document}

%from pagenumber to up
% \phantomsection 
% \hypertarget{toc}{}

 

 

 
   

\section{Закон Кулона. Электрическое поле}


\begin{law}
\textbf{Закон Кулона.} Два тела с зарядами~$q_1$ и~$q_2$, расположенные в вакууме на расстоянии~$r$, взаимодействуют друг с другом с \textbf{силой Кулона}
\begin{equation}
    \label{6Claw_e1}
    F_\text{К}=k\dfrac{q_1q_2}{r^2},
\end{equation}
где~$k$~--- коэффициент пропорциональности (см.~cправочные таблицы). 
\end{law}

На рис.\,\ref{fig:p96} изображены два отрицательных заряда.

    \begin{figure}[H]
    \centering

    \begin{tikzpicture}[font=\small,            line cap=round,                             >={Stealth[inset=0pt,round,length=3mm, angle'=20]},vec/.style={>={Stealth[inset=0pt,round,length=3.5mm, angle'=20]}}]


\shade[ball color=NavyBlue] (0,0) circle (.6);

\path (0,0) ++ (0:3cm) coordinate (s);
\shade[ball color=NavyBlue] (s) circle (.15);


%r
\draw[densely dashed] (0,0) --++ (90:1.5cm);
\draw[densely dashed] (s.center) --++ (90:1.5cm);

\draw[<->] (0,0) ++ (90:1.3cm) --++ (0:3cm) node[midway,above] {$r$};


%vec
\node at (0,0) [circle,fill=red,inner sep=0pt,minimum size=1mm,label=below:] {};
\draw[->,red,thick,vec] (0,0) --++ (180:1.2cm) node[black,above] {$F_\text{К}$};
\node at (s) [circle,fill=red,inner sep=0pt,minimum size=1mm,label=below:] {};
\draw[->,red,thick,vec] (s.center) --++ (0:1.2cm) node[black,above] {$F_\text{К}$};

%m
\node [below,black] at (0,-.6) {$q_\text{1}$};
\path (0,0) ++ (0:3cm) ++ (0,-.15) node[black,below] {$q_\text{2}$};



    \end{tikzpicture}
\caption{Взаимодействие двух отрицательных зарядов}
\label{fig:p96}
\end{figure}

Заряды~$q_1$ и~$q_2$ расположены на некотором расстоянии~$r$ друг от друга\footnote{Для однородно заряженных шарообразных тел~$r$ есть расстояние между их центрами.}. Как и любые два~заряженных тела, эти заряды \emph{взаимодействуют} друг с другом. В данном случае заряд~$q_1$ отталкивается от заряда~$q_2$ с силой~$F_\text{К}$, а заряд~$q_2$~--- от заряда~$q_1$ с такой же силой~$F_\text{К}$. 

Сила Кулона зависит от среды, в которой находятся заряды. Часто заряды оказываются помещенными в так называемый \emph{диэлектрик}~--- вещество, которое не проводят через себя электрические заряды (например, если соединить диэлектриком\footnote{Примерами диэлектриков являются пластик, стекло, масло и~т.\,д.} два металлических шара, несущих разноименные заряды, то перераспределения зарядов между шарами не будет). В диэлектрике формула~\eqref{6Claw_e1} приобретает вид:
\begin{equation*}
    F_\text{К}=k\dfrac{q_1q_2}{\varepsilon r^2},
\end{equation*}
где~$\varepsilon$~--- \emph{диэлектрическая проницаемость} среды (см.~cправочные таблицы). 

Считают, что взаимодействие неподвижных зарядов осуществляется посредством \emph{электрического поля}~--- формы материи, окружающей заряженные тела (поле можно представлять себе как невидимое истечение из тела воображаемой жидкости и~т.\,п.). Сказанное иллюстрируется рис.\,\ref{fig:p97}: поле положительного заряда~А <<отталкивает>> положительный заряд~Б от заряда~А (и наоборот)\footnote{Скорость передачи взаимодействия между зарядами равна примерно~$3\cdot10^8$~м/с.}. 
% заряд~Б действет своим полем на заряд~А).


\tikzfading[name=fade out,
inner color=transparent!0,
outer color=transparent!100]

    \begin{figure}[H]
    \centering

    \begin{tikzpicture}[font=\small,            line cap=round,                             >={Stealth[inset=0pt,round,length=3mm, angle'=20]},vec/.style={>={Stealth[inset=0pt,round,length=3.5mm, angle'=20]}}]

\def\df{4}
\fill [red,path fading=fade out] ({-\df/2},{-\df/2}) rectangle++ (\df,\df);
\shade[ball color=red] (0,0) circle (.6);

\path (0,0) ++ (0:1.5cm) coordinate (s);
\shade[ball color=red] (s) circle (.15);


%vec
\node at (s) [circle,fill=red,inner sep=0pt,minimum size=1mm,label=below:] {};
\draw[->,red,thick,vec] (s.center) --++ (0:1.2cm) node[black,above] {$F_\text{К}$};

%m
\node [below,black] at (0,-.6) {А};
\path (0,0) ++ (0:1.5cm) ++ (0,-.15) node[black,below] {Б};



    \end{tikzpicture}
\caption{Действие поля одного заряда на другой}
\label{fig:p97}
\end{figure}

% Скорость передачи взаимодействия между зарядами равна примерно~$3\cdot10^8$~м/с. 
% (скорость света). 

















 
\end{document}
